
% Preamble
\documentclass[12pt, letterpaper]{article}
\title{Chapter 1 Notes}
\author{Jordan Tucker}
\date{April 4, 2026}
% Packages
\usepackage{amsmath}
\usepackage{parskip}
\usepackage{tabularx}

% Document
\begin{document}
    \maketitle

    \section{Introduction to Waves}\label{sec:introduction-to-waves}
    A \textbf{wave} is any recognizable signal that is transferred from one part of a medium to another with
    a recognizable velocity of propagation.
    A medium is a thing a wave passes through, such as water
    or air.

    A wave is caused by a disturbance that disrupts the medium from its normal resting state.
    Think of a rock hitting water or someone speaking.
    That is what causes the wave to begin propagation.

    When observing a wave, one should always be able to track a signal.
    The \textbf{signal} is the specific, measurable feature of that disturbance that you can track
    as it travels through the medium.

    \subsection{Examples of Waves}\label{subsec:examples-of-waves}
    One good example of a wave is a pebble hitting a pond.
    When a pebble hits a pond, it causes a disturbance.
    This disturbance creates a signal, in this case, crests (waves) in the pond.
    This waves propagate outwards from the center of where the pebble hit with a recognizable
    velocity.

    Another good example is car traffic.
    When a car stops at a light, it creates a ``disturbance'' in traffic.
    All the cars behind that car will stop.
    We can track the end of the line as our signal.
    The more cars that get added to the line, the further back our signal moves.

    \begin{tabularx}{\textwidth}{|l | X | X|}
        \hline
        \textbf{Wave Phenomena} & \textbf{Medium} & \textbf{Signal}\\
        \hline
        Ripples on a pond & Surface layer of water & Extreme water surface height \\
        Compression waves & Elastic bar & Maximum longitudinal displacement\\
        Sound waves & Gas or liquid & Pressure extrema \\
        Shock waves & Gas & Abrupt change in pressure \\
        Traffic waves & Car traffic on a road & Abrupt change in traffic density (hitting the end of the line) \\
        Epizootic waves & A population susceptible to contracting a disease & Extreme number of infectives \\
        \hline
    \end{tabularx}

\end{document}